\documentclass[conference]{IEEEtran}

% ── Packages ─────────────────────────────────────────────────────────
\usepackage{amsmath, amssymb, amsthm}
\usepackage{graphicx}
\usepackage{booktabs}
\usepackage{hyperref}
\usepackage{cite}
\usepackage{array}
\usepackage{xcolor}

\hypersetup{
    colorlinks = true,
    linkcolor  = blue,
    citecolor  = blue,
    urlcolor   = blue
}

% ─────────────────────────────────────────────────────────────────────

\begin{document}

\title{Markowitz Portfolio Optimisation Using a From-Scratch Implementation
of Wolfe's Modified Simplex Method for Quadratic Programming}

\author{
    \IEEEauthorblockN{Debdan Samanta, Musaib Bin Bashir, Srijan Ghosh}
    \IEEEauthorblockA{
        Department of Industrial Engineering\\
        Operations Research II\\
    }
}

\maketitle

% ── Abstract ─────────────────────────────────────────────────────────
\begin{abstract}
This paper presents a complete, ground-up implementation of the Markowitz
Mean-Variance Portfolio Optimisation model, solved entirely without the use
of any external optimisation library. The core engine is built on Wolfe's
Modified Simplex Method, a classical technique from 1959 that converts a
convex Quadratic Program~(QP) into an augmented Linear Program via the
Karush--Kuhn--Tucker~(KKT) first-order optimality conditions, and then
solves the resulting system using a modified Simplex pivot procedure that
enforces complementary slackness at every iteration. The implementation is
applied to a real-world portfolio of nine Nifty50 equities and a Gold ETF,
with live price data fetched from Yahoo Finance. Experimental results
demonstrate that the solver correctly traces the Markowitz Efficient Frontier
and produces allocation weights consistent with theoretical expectations.
\end{abstract}

\begin{IEEEkeywords}
Convex Quadratic Programming, Efficient Frontier, KKT Conditions,
Markowitz Portfolio Theory, Wolfe's Modified Simplex Method.
\end{IEEEkeywords}

% ─────────────────────────────────────────────────────────────────────
\section{Introduction}
% ─────────────────────────────────────────────────────────────────────

One of the most well-known problems in quantitative finance is the Markowitz
portfolio selection problem~\cite{markowitz1952}. The core idea is
surprisingly intuitive: given a set of assets, each with some expected return
and some risk (measured as variance), find the combination of assets that
achieves the lowest possible risk for a desired level of return. This is not
just a finance problem --- it is, at its heart, a constrained optimisation
problem, specifically a Quadratic Program.

Most practical implementations of portfolio optimisers rely on off-the-shelf
solvers such as \texttt{scipy.optimize.minimize} or commercial tools like
CVXPY. The goal of this project, however, was to implement the entire solver
pipeline from scratch using only NumPy, in order to understand what is
actually happening under the hood. This turned out to be significantly more
involved than initially expected.

The solver used here is Wolfe's Modified Simplex Method~\cite{wolfe1959},
published in \textit{Econometrica} in 1959. Wolfe's key insight was that the
KKT conditions of a convex QP form a structured linear system, and that a
specially modified version of the Simplex algorithm --- one that respects the
complementarity constraints at every pivot step --- can solve it exactly. This
paper documents the mathematical foundations, the algorithmic design, and the
implementation decisions made throughout this project.

% ─────────────────────────────────────────────────────────────────────
\section{Problem Formulation}
% ─────────────────────────────────────────────────────────────────────

\subsection{The Markowitz Mean-Variance Model}

Let there be $n$ assets available for investment. The decision variable is the
weight vector $w \in \mathbb{R}^n$, where each $w_i$ represents the percentage
of total capital allocated to asset $i$. The relevant parameters are:

\begin{itemize}
    \item $\Sigma \in \mathbb{R}^{n \times n}$: the annualised covariance matrix
          of asset returns (in units of $\%^2$), which encodes both individual
          asset risk and pairwise correlations.
    \item $r \in \mathbb{R}^n$: the vector of annualised expected returns (in \%).
    \item $T \in \mathbb{R}$: the investor's target minimum annualised return
          (in \%).
\end{itemize}

The portfolio optimisation problem is stated as follows:

\begin{equation}
    \begin{aligned}
        \min_{w} \quad        & w^\top \Sigma\, w                                \\
        \text{subject to}\quad & \mathbf{1}^\top w = 100                          \\
                               & r^\top w \geq 100 \cdot T                        \\
                               & w_i \geq 0, \quad \forall\, i \in \{1,\dots,n\}
    \end{aligned}
    \label{eq:markowitz}
\end{equation}

The objective minimises portfolio variance, which is a direct measure of risk.
The first constraint enforces full capital deployment (no cash position). The
second constraint requires that the portfolio achieves at least the target
return. The third constraint enforces that no asset is shorted (long-only
positions).

Note that $\Sigma \succeq 0$ (positive semi-definite) by construction, which
means the objective is convex. This is an important property because it
guarantees that any locally optimal solution found is also globally optimal.

\subsection{Reformulation as a Canonical QP}

Wolfe's solver requires the problem to be in the following standard form:

\begin{equation}
    \min_{x \geq 0} \;\; \frac{1}{2} x^\top Q x + c^\top x
    \quad \text{s.t.} \quad
    A_{\text{eq}}\, x = b_{\text{eq}}, \;\;
    A_{\text{ub}}\, x \leq b_{\text{ub}}
    \label{eq:canonical_qp}
\end{equation}

The portfolio problem maps to this canonical form by setting:
\begin{equation}
    Q = 2\Sigma, \quad c = \mathbf{0}, \quad
    A_{\text{eq}} = \mathbf{1}^\top, \quad b_{\text{eq}} = 100
    \label{eq:qp_params}
\end{equation}

The return constraint $r^\top w \geq 100T$ is converted to $\leq$ form by
negating both sides:
\begin{equation}
    A_{\text{ub}} = -r^\top, \qquad b_{\text{ub}} = -100T
    \label{eq:return_ineq}
\end{equation}

The factor of 2 in $Q = 2\Sigma$ is included because the solver minimises
$\frac{1}{2} x^\top Q x$, so $\frac{1}{2} w^\top (2\Sigma) w = w^\top \Sigma w$,
which is exactly what we want.

% ─────────────────────────────────────────────────────────────────────
\section{Wolfe's Modified Simplex Method}
% ─────────────────────────────────────────────────────────────────────

\subsection{Motivation}

The standard Simplex algorithm is designed for Linear Programs (LPs) ---
problems where both the objective and constraints are linear. The Markowitz
problem has a quadratic objective, so it cannot be solved directly with
Simplex. Wolfe's contribution was to transform the QP into an equivalent
LP-like system by leveraging the KKT optimality conditions.

\subsection{KKT Conditions}

For a convex QP, the KKT conditions are both necessary and sufficient for
global optimality~\cite{bazaraa2006}. The Lagrangian of the canonical QP
is:

\begin{equation}
\begin{split}
    \mathcal{L}(w, \lambda, \mu, s) &= \frac{1}{2} w^\top Q w + c^\top w \\
    &\quad + \lambda^\top (A_{\text{eq}} w - b_{\text{eq}}) \\
    &\quad + \mu^\top (A_{\text{ub}} w - b_{\text{ub}}) - s^\top w
\end{split}
\label{eq:lagrangian}
\end{equation}

where $\lambda \in \mathbb{R}^p$ are the (free) Lagrange multipliers for the
equality constraints, $\mu \in \mathbb{R}^q_{\geq 0}$ are the multipliers for
the inequality constraints, and $s \in \mathbb{R}^n_{\geq 0}$ are the dual
surplus variables associated with the non-negativity constraint $w \geq 0$.

Taking the gradient with respect to $w$ and setting it to zero gives the
\textit{stationarity condition}:

\begin{equation}
    Qw + c + A_{\text{eq}}^\top \lambda + A_{\text{ub}}^\top \mu - s = 0
    \label{eq:stationarity}
\end{equation}

The \textit{complementary slackness conditions} state:

\begin{equation}
    s_i \cdot w_i = 0 \;\; \forall\, i, \qquad
    \mu_j \cdot t_j = 0 \;\; \forall\, j
    \label{eq:complementarity}
\end{equation}

where $t_j = b_{\text{ub},j} - (A_{\text{ub}} w)_j \geq 0$ is the slack
variable for the $j$-th inequality. These conditions state that for each
complementary pair, at most one of the two variables can be nonzero at
optimality.

\subsection{The Augmented Tableau}

Wolfe's method constructs an augmented LP tableau whose rows encode the full
KKT system. The variables in column order are listed in Table~\ref{tab:vars}.

\begin{table}[h]
    \centering
    \caption{Tableau Variable Layout}
    \label{tab:vars}
    \begin{tabular}{clp{4.5cm}}
        \toprule
        \textbf{Symbol} & \textbf{Count} & \textbf{Meaning} \\
        \midrule
        $w$               & $n$      & Primal decision variables \\
        $s$               & $n$      & KKT surplus; complementary to $w$ \\
        $\lambda^+,\lambda^-$ & $p$ each & Positive/negative parts of the free multiplier $\lambda = \lambda^+ - \lambda^-$ \\
        $\mu$             & $q$      & Inequality multipliers $\geq 0$; complementary to $t$ \\
        $t$               & $q$      & Inequality slack variables; complementary to $\mu$ \\
        \bottomrule
    \end{tabular}
\end{table}

The rows of the tableau encode the full KKT system:

\begin{align}
    &Qw - s + A_{\text{eq}}^\top(\lambda^+ - \lambda^-) + A_{\text{ub}}^\top \mu = -c
    \label{eq:tableau_stat} \\[4pt]
    &A_{\text{eq}}\, w = b_{\text{eq}}
    \label{eq:tableau_eq} \\[4pt]
    &A_{\text{ub}}\, w + t = b_{\text{ub}}
    \label{eq:tableau_ineq}
\end{align}

This system, when solved with $w, s, \mu, t \geq 0$ and the complementarity
conditions satisfied, yields the optimal solution to the original QP.

\subsection{Phase I: Big-M Initialisation}

The Simplex method requires a starting basis --- a set of $m$ basic variables
whose corresponding tableau columns form an identity matrix. For general QPs,
no obvious initial basis exists. Wolfe's method uses the \textit{Big-M
technique} to handle this.

For every row that lacks a natural identity column, an \textit{artificial
variable} $a_i \geq 0$ is introduced. The objective is then augmented with a
large penalty:

\begin{equation}
    \min \quad \frac{1}{2} w^\top Q w + c^\top w + M \sum_i a_i,
    \qquad M = 10^6
    \label{eq:bigm}
\end{equation}

This drives all artificials to zero during optimisation, effectively finding
a feasible point for the KKT system. If any artificial remains in the basis at
a nonzero value at termination, the original problem is infeasible.

In the portfolio problem specifically, since $c = \mathbf{0}$ and the return
constraint has a negative RHS (which gets sign-flipped during tableau
construction), all stationarity rows and all constraint rows require artificial
variables.

\subsection{The Complementarity Restriction}

This is the key modification that distinguishes Wolfe's method from a vanilla
Simplex. The standard pivot rule selects the entering variable as the one with
the most negative reduced cost. Wolfe adds the following restriction:

\begin{quote}
    \textit{Variable $j$ may enter the basis only if its complementary
    partner is not currently in the basis.}
\end{quote}

The complementary pairs are $(w_i, s_i)$ for all $i$, and $(\mu_j, t_j)$ for
all $j$. Without this restriction, the Simplex method might produce a solution
where both $w_i > 0$ and $s_i > 0$ simultaneously, which would violate the
complementary slackness condition~\eqref{eq:complementarity} and would not be
a valid KKT point.

A two-pass selection procedure is used in practice:

\begin{enumerate}
    \item \textbf{Pass 1}: Find the most negative reduced cost among all
          variables whose complementary partner is not in the basis.
    \item \textbf{Pass 2}: If no such variable exists, check whether any
          blocked variable can enter via a degenerate pivot --- specifically,
          one where the complementary partner would simultaneously leave the
          basis (its min-ratio equals zero). This handles degenerate cases
          without getting stuck.
\end{enumerate}

\subsection{Pivot Step}

Once an entering variable is chosen, the leaving variable is determined by the
standard \textit{minimum ratio test}:

\begin{equation}
    \text{leave} = \arg\min_{\;i\;:\;\text{col}[i]\;>\;0}
    \frac{\text{RHS}[i]}{\text{col}[i]}
    \label{eq:min_ratio}
\end{equation}

The pivot operation performs row reduction to restore the basis columns to
standard form, exactly as in standard Simplex.

\subsection{Termination Criteria}

The algorithm terminates under one of three conditions:

\begin{enumerate}
    \item \textbf{Optimal}: All reduced costs are non-negative and no eligible
          entering variable exists.
    \item \textbf{Infeasible}: An artificial variable remains in the basis with
          a positive value, meaning the constraints are contradictory.
    \item \textbf{Unbounded}: No positive entry exists in the entering column,
          meaning the objective can decrease without bound. This should not
          occur for a convex QP with a bounded feasible region.
\end{enumerate}

% ─────────────────────────────────────────────────────────────────────
\section{Implementation}
% ─────────────────────────────────────────────────────────────────────

\subsection{Numerical Stability}

\textbf{Positive Semi-Definiteness.}
Historical covariance matrices estimated from finite data samples can have
small negative eigenvalues due to sampling noise, which would make the QP
non-convex. To handle this, a spectral decomposition is performed and any
negative eigenvalues are clipped to zero:

\begin{equation}
    \Sigma \;\leftarrow\; V \cdot \max(\Lambda,\, 0) \cdot V^\top
    \label{eq:psd_clip}
\end{equation}

where $\Sigma = V \Lambda V^\top$ is the eigendecomposition. This yields
the nearest PSD matrix to the original in the Frobenius norm sense.

\textbf{Symmetry Enforcement.}
Floating-point arithmetic can introduce tiny asymmetries in $\Sigma$ and $Q$.
Both matrices are explicitly symmetrised at multiple points via
$Q \leftarrow \frac{1}{2}(Q + Q^\top)$.

\textbf{Weight Scaling.}
The solver works with weights summing to 100 (percent) rather than 1
(fraction). This improves the numerical conditioning of the tableau by keeping
all matrix entries and RHS values in the range $10^1$--$10^2$, which reduces
the relative effect of floating-point rounding during pivot operations.

\textbf{Tolerance Constants.}
Three key tolerances govern numerical precision and are listed in
Table~\ref{tab:tols}.

\begin{table}[h]
    \centering
    \caption{Numerical Tolerance Constants}
    \label{tab:tols}
    \begin{tabular}{lll}
        \toprule
        \textbf{Constant} & \textbf{Value} & \textbf{Role} \\
        \midrule
        \texttt{\_BIG\_M}     & $10^6$    & Penalty weight for artificial variables \\
        \texttt{\_PIVOT\_TOL} & $10^{-10}$ & Minimum valid pivot element magnitude \\
        \texttt{\_FEAS\_TOL}  & $10^{-8}$  & Zero threshold for artificial variables \\
        \bottomrule
    \end{tabular}
\end{table}

\subsection{Software Architecture}

The implementation is structured as three independent modules with a clean
dependency hierarchy:

\begin{itemize}
    \item \texttt{QuadraticSimplex.py}: A general-purpose QP solver exposing
          a single public function \texttt{solve\_qp\_wolfe()}. It has no
          project-specific dependencies and can be reused for any convex QP.
    \item \texttt{PortfolioOptimizer.py}: Translates the Markowitz problem
          into the canonical QP form, enforces PSD on the covariance matrix,
          and wraps the solver output into structured result objects.
    \item \texttt{Task.py}: The top-level runner that fetches live market data,
          computes statistics, invokes the optimizer, and produces
          visualisations.
\end{itemize}

The dependency direction is strictly one-way:
$\texttt{Task.py} \to \texttt{PortfolioOptimizer.py} \to \texttt{QuadraticSimplex.py}$.

\subsection{Data Pipeline}

Live daily adjusted-close prices are fetched from Yahoo Finance via the
\texttt{yfinance} library for the most recent four months.
\textit{Log returns} are computed as $r_t = \ln(P_t / P_{t-1})$, preferred
over simple returns because they are time-additive across periods and better
approximate a normal distribution for short intervals --- a key assumption of
the Markowitz model. Annualised statistics are obtained by scaling by 252
trading days per year:

\begin{equation}
    \hat{\mu}_{\text{ann}} = 252 \cdot \bar{r}_{\text{daily}},
    \qquad
    \hat{\Sigma}_{\text{ann}} = 252 \cdot \hat{\Sigma}_{\text{daily}}
    \label{eq:annualise}
\end{equation}

Missing data (caused by thin trading or recent IPOs) is handled by
backward-then-forward filling nearest known prices. Any asset with
irrecoverable missing data is dropped from the universe entirely.

% ─────────────────────────────────────────────────────────────────────
\section{Results}
% ─────────────────────────────────────────────────────────────────────

The solver was applied to a universe of nine Nifty50 equities
(RELIANCE, TCS, HDFCBANK, INFY, ICICIBANK, HINDUNILVR, SBIN, BAJFINANCE,
WIPRO) plus the GoldBEES ETF, using four months of daily price history.

For a user-specified target return $T$, the solver produces the
minimum-variance portfolio allocation. The efficient frontier is traced by
solving the QP for 50 evenly-spaced target return values between the minimum
and maximum feasible returns. The output is a three-panel visualisation:

\begin{itemize}
    \item \textbf{Left panel}: The efficient frontier curve in
          $(\sigma_p, \mu_p)$ space, with individual assets plotted at their
          own risk-return coordinates and the optimal portfolio marked with a
          star.
    \item \textbf{Centre panel}: A bar chart of the optimal weight allocation
          for the chosen target $T$.
    \item \textbf{Right panel}: The evolution of each asset's weight along the
          frontier, showing how diversification shifts as the target return
          increases.
\end{itemize}

The solver correctly finds corner solutions (100\% allocation to a single
asset) at extreme target returns, and interior diversified allocations at
intermediate returns, consistent with theoretical expectations
from~\cite{markowitz1952}.

To illustrate the solver's practical output, two specific target annualised
returns were evaluated: $T = 10\%$ and $T = 15\%$. The resulting optimal
portfolio allocations and their associated risk statistics are summarised in
Table~\ref{tab:portfolios}. Both scenarios yield diversified portfolios across
RELIANCE, SBIN, and Gold, with higher risk targets driving more aggressive
allocations.

\begin{table}[h]
    \centering
    \caption{Optimal Portfolios for $T=10\%$ and $T=15\%$ Targets}
    \label{tab:portfolios}
    \begin{tabular}{lrr}
        \toprule
        \textbf{Metric / Asset Weight} & \textbf{$T = 10\%$} & \textbf{$T = 15\%$} \\
        \midrule
        \textbf{Portfolio Return} & 10.00\% & 15.00\% \\
        \textbf{Portfolio Std Dev ($\sigma$)} & 20.88\% & 22.65\% \\
        \textbf{Portfolio Variance ($\sigma^2$)} & 436.17 $\%^2$ & 512.85 $\%^2$ \\
        \textbf{Sharpe Ratio ($r_f=0$)} & 0.4788 & 0.6624 \\
        \midrule
        RELIANCE & 19.40\% & 8.94\% \\
        SBIN & 65.40\% & 73.90\% \\
        Gold & 15.21\% & 17.16\% \\
        \textit{All other assets} & 0.00\% & 0.00\% \\
        \bottomrule
    \end{tabular}
\end{table}

% ─────────────────────────────────────────────────────────────────────
\section{Conclusion}
% ─────────────────────────────────────────────────────────────────────

This project demonstrates that the Markowitz portfolio optimisation problem
can be solved from first principles, without any external QP library, by
implementing Wolfe's Modified Simplex Method directly. The key contributions
of this implementation are: (i) a complete Big-M Phase~I initialisation that
dynamically identifies which rows require artificial variables; (ii) a
two-pass complementarity restriction that correctly handles degenerate pivot
cases; and (iii) a numerically stable problem formulation with PSD enforcement
and percentage-unit scaling.

One limitation worth noting is that the Big-M method, while straightforward,
can sometimes suffer from numerical issues when $M$ is not large enough to
force artificials out, or when it interacts poorly with the problem's natural
scaling. A more robust alternative would be a two-phase Simplex that separates
Phase~I (feasibility) from Phase~II (optimisation) entirely. That said, for
the problem sizes encountered in this project (10 assets, 2 constraints), the
Big-M approach performs reliably.

% ─────────────────────────────────────────────────────────────────────
\begin{thebibliography}{9}

\bibitem{markowitz1952}
H.~Markowitz,
``Portfolio Selection,''
\textit{The Journal of Finance},
vol.~7, no.~1, pp.~77--91, Mar.\ 1952.

\bibitem{wolfe1959}
P.~Wolfe,
``The Simplex Method for Quadratic Programming,''
\textit{Econometrica},
vol.~27, no.~3, pp.~382--398, Jul.\ 1959.

\bibitem{bazaraa2006}
M.~S.~Bazaraa, H.~D.~Sherali, and C.~M.~Shetty,
\textit{Nonlinear Programming: Theory and Algorithms},
3rd~ed.\ Hoboken, NJ, USA: Wiley, 2006, ch.~4.

\bibitem{chong2013}
E.~K.~P.~Chong and S.~H.~Zak,
\textit{An Introduction to Optimization},
4th~ed.\ Hoboken, NJ, USA: Wiley, 2013, ch.~17.

\bibitem{yfinance}
R.~Aroussi,
``yfinance: Download market data from Yahoo! Finance,''
Python Package Index (PyPI).
[Online]. Available: \url{https://pypi.org/project/yfinance/}

\end{thebibliography}

\end{document}
